
\sourcenote{Genre: admissible rule (tier 2). The family machinery is
stated in the \rulename{LiRPA} section's family template. This rule is
premise-free: $\Phi$ is arbitrary and never consulted (the
\rulename{Norm} convention).}

\begin{brecipe}[the \rulename{ReLU-NN} instance --- the full unpacking]
\label{rec:relunn-instance}
The node:
\[
\inferrule*[Right=\rulename{ReLU-NN}]{
  \ReLU_{\{h\}\to\{y\}} \in H
}{
  \Phi \vdash y \geq 0
}
\qquad\rightsquigarrow\qquad
\text{case analysis on }
(h \leq 0 \wedge y = 0) \vee (0 \leq h \wedge y = h),
\]
and both branches discharge into catalogue nodes --- the family's
cleanest unpacking. \emph{Inactive branch} (context carries $y = 0$):
\begin{enumerate}
\item $\vdash 0 \geq 0$ --- \rulename{PConstExact}, the closed exact
  ground fact (before that rule arrived, this leaf had no derivation at
  all --- the calculus could not state a ground fact);
\item $\Phi_{-} \vdash y = 0$ --- \rulename{Ax} + $\wedge$E from the
  branch context; then $\Phi_{-} \vdash 0 = y$ by \rulename{EqSym};
\item $\Phi_{-} \vdash y \geq 0$ --- \rulename{EqElim} along $0 = y$,
  transporting $\Psi[a \coloneqq 0] = (0 \geq 0)$ (step 1, weakened)
  to $\Psi[a \coloneqq y]$, where $\Psi = (a \geq 0)$.
\end{enumerate}
\emph{Active branch} (context carries $0 \leq h$ and $y = h$): the
context's sign atom is re-oriented to $h \geq 0$ by the polarity
plumbing walked in full in the \rulename{ReLU-Above} section, and
\rulename{EqElim} along $h = y$ (an equality between two wires of the
node's one shared shape) transports it to $y \geq 0$. $\vee$E in its
context-carrying form closes the case analysis, and \rulename{Cut}
stitches the sign-split fact into place.
\end{brecipe}

\begin{bremark}[the shape scope of the displayed derivation]
\label{rem:relunn-shapescope}
The realized \rulename{EqElim} carries a broadcast-shape side
condition (its section documents the counterexample that forced it):
both sides of the transported equality must share a broadcast shape.
Step 3's equality $0 = y$ has a rank-$0$ constant on one side, so the
displayed transport type-checks where the output wire is rank-$0$ ---
at tensor-shaped instances that step is stopped by the gate, and the family's
realized warrant does not route through the derivation at all: it is
the direct lemma \texttt{sound\_reluNN}, shape-general, via the
linkage extraction. (Wire-to-wire transports such as the active
branch's $h = y$ are unaffected: the linkage forces one shared shape.)
The derivation above is thus the honest admissibility story at scalar
instances, and the reduction-to-verified-machinery story --- the
\rulename{PConstExact} pattern, stated in the family template --- is
what covers every shape.
\end{bremark}

\begin{bprop}[semantic warrant]\label{prop:relunn-sound}
If the linkage holds, the sequent concluding $y \geq 0$ is
semantically valid over any premise formula $\Phi$.
\end{bprop}
\begin{proof}
Fix an admissible valuation $\sigma$ ($\sigma \models \Phi$ is taken
but never consulted --- the rule is premise-free). The linkage
extraction (\texttt{reluLinked\_facts}) gives the output tensor and,
per entry, either $y = 0$ or $y = p$ with $0 \leq p$; in both cases
$0 \leq y$. The conclusion atom compares the output tensor entrywise
against the rank-$0$ constant $0$. That is \texttt{sound\_reluNN},
step for step.
\end{proof}

\begin{bremark}[the only hypothesis]
The linkage is the rule's sole side condition, and dropping it is the
whole failure mode: an arbitrary wire asserted nonnegative against an
arbitrary premise. Its positive-size gate additionally keeps the
vacuous corner unciteable --- a size-zero node satisfies every
elementwise atom trivially, and a fact no valuation can falsify must
not enter a derivation (the family template's hygiene remark).
\end{bremark}

\begin{anchors}
Recipe \ref{rec:relunn-instance} warrant & \texttt{sound\_reluNN}\\
family linkage extraction & \texttt{reluLinked\_facts} (see the \rulename{LiRPA} section)\\
ground-fact leaf (step 1) & \texttt{sound\_pConstExact} (the \rulename{PConstExact} section)\\
payload structure & \texttt{PReLUNNData} (two wires, linkage --- no premise field)\\
\end{anchors}
